\documentclass[11pt,letterpaper]{article}
\usepackage{cornell-notes}
\usepackage{needspace}

\newcommand{\CornellDocumentTitle}{Chapter 47: data encryption}
\title{\CornellDocumentTitle}
\author{}
\date{}
\setDocTitle{\CornellDocumentTitle}
\setDocSubtitle{Cybersecurity Cornell Notes}
\setCornellCollection{Cybersecurity Reference Notes}
\setCornellUnitType{Chapter}
\setCornellUnitNumber{47}
\setCornellUnitTitle{data encryption}
\setCornellCompanionLabel{Cybersecurity study companion}

\begin{document}
\maketitle
\begin{CornellOverviewBox}{Chapter Orientation}
\textbf{Chapter orientation.} This chapter examines data encryption within the broader domain of encryption technology. Its central problem is how to reason about message, mod, key, and group while keeping security objectives, operating assumptions, evidence, and recovery connected. A practitioner should approach the material through security properties, algorithms, keys, protocols, validation, trust assumptions, and failure through misuse. Useful prerequisites include basic risk, threat, control, and systems concepts.
\end{CornellOverviewBox}
\section*{Learning Objectives}
\begin{itemize}
\item Explain the security purpose and scope of Data Encryption.
\item Identify the assets, actors, assumptions, and trust boundaries associated with message, mod, and key.
\item Distinguish Chapter orientation from Need For Cryptography and explain where their responsibilities intersect.
\item Analyze how failures in message, mod, and key can affect confidentiality, integrity, availability, privacy, or safety.
\item Evaluate relevant controls through security properties, algorithms, keys, protocols, validation, trust assumptions, and failure through misuse.
\item Audit an implementation using algorithm and mode inventories, key provenance, certificate paths, protocol traces, rotation records, and test vectors.
\item Prioritize remediation according to exploitability, consequence, control strength, and recovery readiness.
\item Verify that corrective actions change observable risk rather than documentation alone.
\end{itemize}
\section*{Cornell Cue-and-Notes}
\Needspace{8\baselineskip}
\subsection*{Chapter orientation}
\begin{longtable}{>{\RaggedRight\arraybackslash}p{0.29\textwidth}|>{\RaggedRight\arraybackslash}p{0.65\textwidth}}
\CornellHeader
\CornellNoteRow{What security problem does Chapter orientation address?}{\textbf{Core idea.} Treat Chapter orientation as a structured part of Data Encryption, with particular attention to activities, design, users, and violation allowed. \textbf{Analytical lens.} This topic establishes a component of the chapter's argument; connect its definitions and mechanisms to a testable security objective. For activities, design, and users, the security claim is only as strong as key custody, randomness, parameter choice, validation, and protocol context. \textbf{Security reasoning.} Identify what must remain protected, who or what can alter the expected state, which assumptions cross a trust boundary, and how failure changes confidentiality, integrity, availability, privacy, or safety. \textbf{Application.} The practitioner should verify the complete cryptographic construction and lifecycle rather than the algorithm name alone. \textbf{Evidence.} Seek algorithm and mode inventories, key provenance, certificate paths, protocol traces, rotation records, and test vectors.}
\end{longtable}
\begin{CornellSummaryBox}{Topic Summary}
Chapter orientation is best remembered as the connection between activities, design, users, and violation allowed and the chapter's larger control objective. A defensible review states the assumption, expected behavior, evidence, failure consequence, and required response.
\end{CornellSummaryBox}
\Needspace{8\baselineskip}
\subsection*{Need For Cryptography}
\begin{longtable}{>{\RaggedRight\arraybackslash}p{0.29\textwidth}|>{\RaggedRight\arraybackslash}p{0.65\textwidth}}
\CornellHeader
\CornellNoteRow{Which assumptions matter in Need For Cryptography?}{\textbf{Core idea.} Treat Need For Cryptography as a structured part of Data Encryption, with particular attention to bob, message, alice, and quantum. \textbf{Analytical lens.} This topic is cryptography-centered: state the intended property and validate algorithms, parameters, keys, protocol binding, and lifecycle controls. For bob, message, and alice, the security claim is only as strong as key custody, randomness, parameter choice, validation, and protocol context. \textbf{Security reasoning.} Identify what must remain protected, who or what can alter the expected state, which assumptions cross a trust boundary, and how failure changes confidentiality, integrity, availability, privacy, or safety. \textbf{Application.} The practitioner should verify the complete cryptographic construction and lifecycle rather than the algorithm name alone. \textbf{Evidence.} Seek algorithm and mode inventories, key provenance, certificate paths, protocol traces, rotation records, and test vectors. Related subtopics include Authentication, Confidentiality, Integrity, and Nonrepudiation.}
\end{longtable}
\begin{CornellSummaryBox}{Topic Summary}
Need For Cryptography is best remembered as the connection between bob, message, alice, and quantum and the chapter's larger control objective. A defensible review states the assumption, expected behavior, evidence, failure consequence, and required response.
\end{CornellSummaryBox}
\Needspace{8\baselineskip}
\subsection*{Mathematical Prelude To Cryptography}
\begin{longtable}{>{\RaggedRight\arraybackslash}p{0.29\textwidth}|>{\RaggedRight\arraybackslash}p{0.65\textwidth}}
\CornellHeader
\CornellNoteRow{How should a reviewer evaluate Mathematical Prelude To Cryptography?}{\textbf{Core idea.} Treat Mathematical Prelude To Cryptography as a structured part of Data Encryption, with particular attention to message, function, polynomial, and time. \textbf{Analytical lens.} This topic is cryptography-centered: state the intended property and validate algorithms, parameters, keys, protocol binding, and lifecycle controls. For message, function, and polynomial, the security claim is only as strong as key custody, randomness, parameter choice, validation, and protocol context. \textbf{Security reasoning.} Identify what must remain protected, who or what can alter the expected state, which assumptions cross a trust boundary, and how failure changes confidentiality, integrity, availability, privacy, or safety. \textbf{Application.} The practitioner should verify the complete cryptographic construction and lifecycle rather than the algorithm name alone. \textbf{Evidence.} Seek algorithm and mode inventories, key provenance, certificate paths, protocol traces, rotation records, and test vectors. Related subtopics include Complexity, and Mapping or Function.}
\end{longtable}
\begin{CornellSummaryBox}{Topic Summary}
Mathematical Prelude To Cryptography is best remembered as the connection between message, function, polynomial, and time and the chapter's larger control objective. A defensible review states the assumption, expected behavior, evidence, failure consequence, and required response.
\end{CornellSummaryBox}
\Needspace{8\baselineskip}
\subsection*{Classical Cryptography}
\begin{longtable}{>{\RaggedRight\arraybackslash}p{0.29\textwidth}|>{\RaggedRight\arraybackslash}p{0.65\textwidth}}
\CornellHeader
\CornellNoteRow{Why can Classical Cryptography fail in practice?}{\textbf{Core idea.} Treat Classical Cryptography as a structured part of Data Encryption, with particular attention to mod, inverse, integers, and set. \textbf{Analytical lens.} This topic is cryptography-centered: state the intended property and validate algorithms, parameters, keys, protocol binding, and lifecycle controls. For mod, inverse, and integers, the security claim is only as strong as key custody, randomness, parameter choice, validation, and protocol context. \textbf{Security reasoning.} Identify what must remain protected, who or what can alter the expected state, which assumptions cross a trust boundary, and how failure changes confidentiality, integrity, availability, privacy, or safety. \textbf{Application.} The practitioner should verify the complete cryptographic construction and lifecycle rather than the algorithm name alone. \textbf{Evidence.} Seek algorithm and mode inventories, key provenance, certificate paths, protocol traces, rotation records, and test vectors. Related subtopics include Probability, Modular Arithmetic, The Euclidean Algorithm, and Congruence.}
\end{longtable}
\begin{CornellSummaryBox}{Topic Summary}
Classical Cryptography is best remembered as the connection between mod, inverse, integers, and set and the chapter's larger control objective. A defensible review states the assumption, expected behavior, evidence, failure consequence, and required response.
\end{CornellSummaryBox}
\Needspace{8\baselineskip}
\subsection*{Modern Symmetric Ciphers}
\begin{longtable}{>{\RaggedRight\arraybackslash}p{0.29\textwidth}|>{\RaggedRight\arraybackslash}p{0.65\textwidth}}
\CornellHeader
\CornellNoteRow{What security problem does Modern Symmetric Ciphers address?}{\textbf{Core idea.} Treat Modern Symmetric Ciphers as a structured part of Data Encryption, with particular attention to key, algorithm, message, and cipher. \textbf{Analytical lens.} This topic is cryptography-centered: state the intended property and validate algorithms, parameters, keys, protocol binding, and lifecycle controls. For key, algorithm, and message, the security claim is only as strong as key custody, randomness, parameter choice, validation, and protocol context. \textbf{Security reasoning.} Identify what must remain protected, who or what can alter the expected state, which assumptions cross a trust boundary, and how failure changes confidentiality, integrity, availability, privacy, or safety. \textbf{Application.} The practitioner should verify the complete cryptographic construction and lifecycle rather than the algorithm name alone. \textbf{Evidence.} Seek algorithm and mode inventories, key provenance, certificate paths, protocol traces, rotation records, and test vectors. Related subtopics include Substitution Box, and Permutation Boxes.}
\end{longtable}
\begin{CornellSummaryBox}{Topic Summary}
Modern Symmetric Ciphers is best remembered as the connection between key, algorithm, message, and cipher and the chapter's larger control objective. A defensible review states the assumption, expected behavior, evidence, failure consequence, and required response.
\end{CornellSummaryBox}
\Needspace{8\baselineskip}
\subsection*{Algebraic Structure}
\begin{longtable}{>{\RaggedRight\arraybackslash}p{0.29\textwidth}|>{\RaggedRight\arraybackslash}p{0.65\textwidth}}
\CornellHeader
\CornellNoteRow{Which assumptions matter in Algebraic Structure?}{\textbf{Core idea.} Treat Algebraic Structure as a structured part of Data Encryption, with particular attention to field, polynomial, group, and finite. \textbf{Analytical lens.} This topic establishes a component of the chapter's argument; connect its definitions and mechanisms to a testable security objective. For field, polynomial, and group, the security claim is only as strong as key custody, randomness, parameter choice, validation, and protocol context. \textbf{Security reasoning.} Identify what must remain protected, who or what can alter the expected state, which assumptions cross a trust boundary, and how failure changes confidentiality, integrity, availability, privacy, or safety. \textbf{Application.} The practitioner should verify the complete cryptographic construction and lifecycle rather than the algorithm name alone. \textbf{Evidence.} Seek algorithm and mode inventories, key provenance, certificate paths, protocol traces, rotation records, and test vectors. Related subtopics include Definition Group, Product Ciphers, Definitions of Finite and Infinite Groups, and Rings.}
\end{longtable}
\begin{CornellSummaryBox}{Topic Summary}
Algebraic Structure is best remembered as the connection between field, polynomial, group, and finite and the chapter's larger control objective. A defensible review states the assumption, expected behavior, evidence, failure consequence, and required response.
\end{CornellSummaryBox}
\Needspace{8\baselineskip}
\subsection*{The Internal Functions Of Rijndael In Advanced Encryption Standard Implementation}
\begin{longtable}{>{\RaggedRight\arraybackslash}p{0.29\textwidth}|>{\RaggedRight\arraybackslash}p{0.65\textwidth}}
\CornellHeader
\CornellNoteRow{How should a reviewer evaluate The Internal Functions Of Rijndael In Advanced Encryption Standard Implementation?}{\textbf{Core idea.} Treat The Internal Functions Of Rijndael In Advanced Encryption Standard Implementation as a structured part of Data Encryption, with particular attention to byte, key, state, and round. \textbf{Analytical lens.} This topic is cryptography-centered: state the intended property and validate algorithms, parameters, keys, protocol binding, and lifecycle controls. For byte, key, and state, the security claim is only as strong as key custody, randomness, parameter choice, validation, and protocol context. \textbf{Security reasoning.} Identify what must remain protected, who or what can alter the expected state, which assumptions cross a trust boundary, and how failure changes confidentiality, integrity, availability, privacy, or safety. \textbf{Application.} The practitioner should verify the complete cryptographic construction and lifecycle rather than the algorithm name alone. \textbf{Evidence.} Seek algorithm and mode inventories, key provenance, certificate paths, protocol traces, rotation records, and test vectors. Related subtopics include Mathematical Preliminaries, State, ShiftRows, and Subkey Addition.}
\end{longtable}
\begin{CornellSummaryBox}{Topic Summary}
The Internal Functions Of Rijndael In Advanced Encryption Standard Implementation is best remembered as the connection between byte, key, state, and round and the chapter's larger control objective. A defensible review states the assumption, expected behavior, evidence, failure consequence, and required response.
\end{CornellSummaryBox}
\Needspace{8\baselineskip}
\subsection*{Use Of Modern Block Ciphers}
\begin{longtable}{>{\RaggedRight\arraybackslash}p{0.29\textwidth}|>{\RaggedRight\arraybackslash}p{0.65\textwidth}}
\CornellHeader
\CornellNoteRow{Why can Use Of Modern Block Ciphers fail in practice?}{\textbf{Core idea.} Treat Use Of Modern Block Ciphers as a structured part of Data Encryption, with particular attention to key, number, blocks, and secret. \textbf{Analytical lens.} This topic is cryptography-centered: state the intended property and validate algorithms, parameters, keys, protocol binding, and lifecycle controls. For key, number, and blocks, the security claim is only as strong as key custody, randomness, parameter choice, validation, and protocol context. \textbf{Security reasoning.} Identify what must remain protected, who or what can alter the expected state, which assumptions cross a trust boundary, and how failure changes confidentiality, integrity, availability, privacy, or safety. \textbf{Application.} The practitioner should verify the complete cryptographic construction and lifecycle rather than the algorithm name alone. \textbf{Evidence.} Seek algorithm and mode inventories, key provenance, certificate paths, protocol traces, rotation records, and test vectors. Related subtopics include Review: Number Theory, The Electronic Code Book, Cipher-Block Chaining, and Coprimes.}
\end{longtable}
\begin{CornellSummaryBox}{Topic Summary}
Use Of Modern Block Ciphers is best remembered as the connection between key, number, blocks, and secret and the chapter's larger control objective. A defensible review states the assumption, expected behavior, evidence, failure consequence, and required response.
\end{CornellSummaryBox}
\Needspace{8\baselineskip}
\subsection*{Public-Key Cryptography}
\begin{longtable}{>{\RaggedRight\arraybackslash}p{0.29\textwidth}|>{\RaggedRight\arraybackslash}p{0.65\textwidth}}
\CornellHeader
\CornellNoteRow{What security problem does Public-Key Cryptography address?}{\textbf{Core idea.} Treat Public-Key Cryptography as a structured part of Data Encryption, with particular attention to order, group, number, and prime. \textbf{Analytical lens.} This topic is cryptography-centered: state the intended property and validate algorithms, parameters, keys, protocol binding, and lifecycle controls. For order, group, and number, the security claim is only as strong as key custody, randomness, parameter choice, validation, and protocol context. \textbf{Security reasoning.} Identify what must remain protected, who or what can alter the expected state, which assumptions cross a trust boundary, and how failure changes confidentiality, integrity, availability, privacy, or safety. \textbf{Application.} The practitioner should verify the complete cryptographic construction and lifecycle rather than the algorithm name alone. \textbf{Evidence.} Seek algorithm and mode inventories, key provenance, certificate paths, protocol traces, rotation records, and test vectors. Related subtopics include Factoring, Fermat's Little Theorem, Discrete Logarithm, and Primitive Roots.}
\end{longtable}
\begin{CornellSummaryBox}{Topic Summary}
Public-Key Cryptography is best remembered as the connection between order, group, number, and prime and the chapter's larger control objective. A defensible review states the assumption, expected behavior, evidence, failure consequence, and required response.
\end{CornellSummaryBox}
\Needspace{8\baselineskip}
\subsection*{Cryptanalysis Of Rivesteshamireadleman}
\begin{longtable}{>{\RaggedRight\arraybackslash}p{0.29\textwidth}|>{\RaggedRight\arraybackslash}p{0.65\textwidth}}
\CornellHeader
\CornellNoteRow{Which assumptions matter in Cryptanalysis Of Rivesteshamireadleman?}{\textbf{Core idea.} Treat Cryptanalysis Of Rivesteshamireadleman as a structured part of Data Encryption, with particular attention to mod, roots, primitive, and group. \textbf{Analytical lens.} This topic is evidence-centered: define observable signals, collection points, analytic limits, escalation criteria, and preservation needs. For mod, roots, and primitive, the security claim is only as strong as key custody, randomness, parameter choice, validation, and protocol context. \textbf{Security reasoning.} Identify what must remain protected, who or what can alter the expected state, which assumptions cross a trust boundary, and how failure changes confidentiality, integrity, availability, privacy, or safety. \textbf{Application.} The practitioner should verify the complete cryptographic construction and lifecycle rather than the algorithm name alone. \textbf{Evidence.} Seek algorithm and mode inventories, key provenance, certificate paths, protocol traces, rotation records, and test vectors. Related subtopics include Factorization Attack, and Discrete Logarithm Problem.}
\end{longtable}
\begin{CornellSummaryBox}{Topic Summary}
Cryptanalysis Of Rivesteshamireadleman is best remembered as the connection between mod, roots, primitive, and group and the chapter's larger control objective. A defensible review states the assumption, expected behavior, evidence, failure consequence, and required response.
\end{CornellSummaryBox}
\Needspace{8\baselineskip}
\subsection*{Diffieehellman Algorithm}
\begin{longtable}{>{\RaggedRight\arraybackslash}p{0.29\textwidth}|>{\RaggedRight\arraybackslash}p{0.65\textwidth}}
\CornellHeader
\CornellNoteRow{How should a reviewer evaluate Diffieehellman Algorithm?}{\textbf{Core idea.} Treat Diffieehellman Algorithm as a structured part of Data Encryption, with particular attention to mod, key, secret, and public. \textbf{Analytical lens.} This topic establishes a component of the chapter's argument; connect its definitions and mechanisms to a testable security objective. For mod, key, and secret, the security claim is only as strong as key custody, randomness, parameter choice, validation, and protocol context. \textbf{Security reasoning.} Identify what must remain protected, who or what can alter the expected state, which assumptions cross a trust boundary, and how failure changes confidentiality, integrity, availability, privacy, or safety. \textbf{Application.} The practitioner should verify the complete cryptographic construction and lifecycle rather than the algorithm name alone. \textbf{Evidence.} Seek algorithm and mode inventories, key provenance, certificate paths, protocol traces, rotation records, and test vectors. Related subtopics include Diffie-Hellman Problem.}
\end{longtable}
\begin{CornellSummaryBox}{Topic Summary}
Diffieehellman Algorithm is best remembered as the connection between mod, key, secret, and public and the chapter's larger control objective. A defensible review states the assumption, expected behavior, evidence, failure consequence, and required response.
\end{CornellSummaryBox}
\Needspace{8\baselineskip}
\subsection*{Elliptic Curve Cryptosystems}
\begin{longtable}{>{\RaggedRight\arraybackslash}p{0.29\textwidth}|>{\RaggedRight\arraybackslash}p{0.65\textwidth}}
\CornellHeader
\CornellNoteRow{Why can Elliptic Curve Cryptosystems fail in practice?}{\textbf{Core idea.} Treat Elliptic Curve Cryptosystems as a structured part of Data Encryption, with particular attention to mod, points, defined, and z23. \textbf{Analytical lens.} This topic is cryptography-centered: state the intended property and validate algorithms, parameters, keys, protocol binding, and lifecycle controls. For mod, points, and defined, the security claim is only as strong as key custody, randomness, parameter choice, validation, and protocol context. \textbf{Security reasoning.} Identify what must remain protected, who or what can alter the expected state, which assumptions cross a trust boundary, and how failure changes confidentiality, integrity, availability, privacy, or safety. \textbf{Application.} The practitioner should verify the complete cryptographic construction and lifecycle rather than the algorithm name alone. \textbf{Evidence.} Seek algorithm and mode inventories, key provenance, certificate paths, protocol traces, rotation records, and test vectors. Related subtopics include An Example, Addition Formula, and Example of Elliptic Curve Addition.}
\end{longtable}
\begin{CornellSummaryBox}{Topic Summary}
Elliptic Curve Cryptosystems is best remembered as the connection between mod, points, defined, and z23 and the chapter's larger control objective. A defensible review states the assumption, expected behavior, evidence, failure consequence, and required response.
\end{CornellSummaryBox}
\Needspace{8\baselineskip}
\subsection*{Message Integrity And Authentication}
\begin{longtable}{>{\RaggedRight\arraybackslash}p{0.29\textwidth}|>{\RaggedRight\arraybackslash}p{0.65\textwidth}}
\CornellHeader
\CornellNoteRow{What security problem does Message Integrity And Authentication address?}{\textbf{Core idea.} Treat Message Integrity And Authentication as a structured part of Data Encryption, with particular attention to message, alice, hash, and bob. \textbf{Analytical lens.} This topic is identity-centered: distinguish identification, authentication, authorization, session state, privilege, and accountability. For message, alice, and hash, the security claim is only as strong as key custody, randomness, parameter choice, validation, and protocol context. \textbf{Security reasoning.} Identify what must remain protected, who or what can alter the expected state, which assumptions cross a trust boundary, and how failure changes confidentiality, integrity, availability, privacy, or safety. \textbf{Application.} The practitioner should verify the complete cryptographic construction and lifecycle rather than the algorithm name alone. \textbf{Evidence.} Seek algorithm and mode inventories, key provenance, certificate paths, protocol traces, rotation records, and test vectors. Related subtopics include Cryptographic Hash Functions, Preimage Resistance, Elliptic Curve Security, and Second Preimage Resistance (Weak Collision.}
\end{longtable}
\begin{CornellSummaryBox}{Topic Summary}
Message Integrity And Authentication is best remembered as the connection between message, alice, hash, and bob and the chapter's larger control objective. A defensible review states the assumption, expected behavior, evidence, failure consequence, and required response.
\end{CornellSummaryBox}
\Needspace{8\baselineskip}
\subsection*{Triple Data Encryption Algorithm Block Cipher}
\begin{longtable}{>{\RaggedRight\arraybackslash}p{0.29\textwidth}|>{\RaggedRight\arraybackslash}p{0.65\textwidth}}
\CornellHeader
\CornellNoteRow{Which assumptions matter in Triple Data Encryption Algorithm Block Cipher?}{\textbf{Core idea.} Treat Triple Data Encryption Algorithm Block Cipher as a structured part of Data Encryption, with particular attention to tdea, cryptographic, physical, and storage. \textbf{Analytical lens.} This topic is cryptography-centered: state the intended property and validate algorithms, parameters, keys, protocol binding, and lifecycle controls. For tdea, cryptographic, and physical, the security claim is only as strong as key custody, randomness, parameter choice, validation, and protocol context. \textbf{Security reasoning.} Identify what must remain protected, who or what can alter the expected state, which assumptions cross a trust boundary, and how failure changes confidentiality, integrity, availability, privacy, or safety. \textbf{Application.} The practitioner should verify the complete cryptographic construction and lifecycle rather than the algorithm name alone. \textbf{Evidence.} Seek algorithm and mode inventories, key provenance, certificate paths, protocol traces, rotation records, and test vectors.}
\end{longtable}
\begin{CornellSummaryBox}{Topic Summary}
Triple Data Encryption Algorithm Block Cipher is best remembered as the connection between tdea, cryptographic, physical, and storage and the chapter's larger control objective. A defensible review states the assumption, expected behavior, evidence, failure consequence, and required response.
\end{CornellSummaryBox}
\begin{CornellWarningBox}{Historical Context}
Some products, protocols, examples, threat observations, or legal references in this chapter are historically bounded. Retain the underlying threat model and control objective, but verify present applicability before treating a named implementation as current guidance.
\end{CornellWarningBox}
\begin{CornellOverviewBox}{Defensive Guidance}
Apply the chapter by making assets, authority, trust, expected behavior, and failure response explicit. In practice, verify the complete cryptographic construction and lifecycle rather than the algorithm name alone. A control is not fully supported until its operation, monitoring, ownership, and recovery behavior are demonstrated with evidence.
\end{CornellOverviewBox}
\section*{Threat-and-Control Map}
\begingroup\scriptsize\setlength{\tabcolsep}{2pt}
\begin{longtable}{>{\RaggedRight\arraybackslash}p{0.135\textwidth}>{\RaggedRight\arraybackslash}p{0.145\textwidth}>{\RaggedRight\arraybackslash}p{0.155\textwidth}>{\RaggedRight\arraybackslash}p{0.145\textwidth}>{\RaggedRight\arraybackslash}p{0.16\textwidth}>{\RaggedRight\arraybackslash}p{0.17\textwidth}}
\toprule\rowcolor{Navy}\color{white}\textbf{Asset} & \color{white}\textbf{Threat / Failure} & \color{white}\textbf{Vulnerability} & \color{white}\textbf{Impact} & \color{white}\textbf{Detection / Evidence} & \color{white}\textbf{Control}\\\midrule
\endfirsthead
\toprule\rowcolor{Navy}\color{white}\textbf{Asset} & \color{white}\textbf{Threat / Failure} & \color{white}\textbf{Vulnerability} & \color{white}\textbf{Impact} & \color{white}\textbf{Detection / Evidence} & \color{white}\textbf{Control}\\\midrule
\endhead
Protected data, cryptographic keys, and trust relationships & Disclosure, forgery, replay, downgrade, or trust substitution & Weak parameters, protocol misuse, poor key management, or failed validation & Broken confidentiality, authenticity, integrity, or nonrepudiation goals & Configuration review, protocol analysis, certificate validation, and lifecycle evidence & Approved constructions, protected keys, sound randomness, validation, and rotation\\\midrule
Assumptions surrounding message & Misuse or unexpected state transition & Trust in mod is not validated & A local weakness becomes a broader compromise path & Review evidence for key and test negative paths & Make trust explicit, constrain authority, and fail safely\\\midrule
Recovery capability and decision evidence & Control failure remains unnoticed or cannot be reversed & Monitoring, ownership, or restoration has not been exercised & Longer exposure and slower, less reliable recovery & Alert-to-response exercises and restoration tests & Assigned ownership, actionable telemetry, preserved evidence, and rehearsed recovery\\\midrule
\bottomrule
\end{longtable}
\endgroup
\section*{Practical Security Checklist}
\begin{itemize}[label=$\square$]
\item Define the in-scope assets, users, services, data, and operational dependencies for Data Encryption.
\item Document the expected behavior and security objective of message.
\item Map identities, privileges, trust boundaries, and externally controlled inputs associated with message and mod.
\item Record the threat or failure being addressed, its prerequisites, likely path, and credible consequences.
\item Inspect algorithm and mode inventories, key provenance, certificate paths, protocol traces, rotation records, and test vectors.
\item Test both the intended path and negative paths, including invalid state, partial failure, excessive privilege, and unavailable dependencies.
\item Confirm preventive controls are complemented by actionable detection and a named response owner.
\item Exercise containment, restoration, and key; record actual recovery behavior and unresolved dependencies.
\item Validate exceptions, compensating controls, expiration dates, and accountable approval.
\item Preserve reproducible evidence and tie each finding to affected assets, risk, remediation, and verification criteria.
\end{itemize}
\begin{CornellSummaryBox}{Chapter Synthesis}
The chapter's principal lesson is that data encryption must be evaluated as an operating security system rather than an isolated product or statement. The most important failure patterns involve message, mod, key, and group, especially when ownership, boundaries, telemetry, or recovery remain implicit. A reviewer should connect architecture and behavior to algorithm and mode inventories, key provenance, certificate paths, protocol traces, rotation records, and test vectors, prioritize findings by reachable consequence and control independence, and verify remediation through observable change. Within Encryption Technology, this chapter contributes a reusable method: state the objective, expose assumptions, test failure, preserve evidence, and learn from operation.
\end{CornellSummaryBox}
\section*{Self-Test Questions}
\begin{enumerate}
\item What is the central security objective of Data Encryption?
\item How does Chapter orientation contribute to the chapter's broader security argument?
\item Distinguish Need For Cryptography from Mathematical Prelude To Cryptography.
\item Which trust assumptions should be made explicit when evaluating message?
\item How could a weakness involving mod affect confidentiality, integrity, availability, privacy, or safety?
\item What evidence would demonstrate that controls surrounding key operate as intended?
\item Why should prevention, detection, response, and recovery be evaluated together in this domain?
\item Scenario: A system uses a recognized algorithm but leaves key generation, validation, and rotation assumptions implicit. What should the reviewer examine first, and why?
\item Scenario: A control associated with group passes a configuration check but fails during an operational exercise. How should the finding be interpreted and prioritized?
\item How should a practitioner distinguish enduring principles in this chapter from time-sensitive tools, examples, or implementation details?
\end{enumerate}
\Needspace{8\baselineskip}
\section*{Answer Key}
\begin{enumerate}
\item The objective is to protect the relevant assets by understanding security properties, algorithms, keys, protocols, validation, trust assumptions, and failure through misuse and by connecting controls to measurable risk reduction.
\item Chapter orientation provides one part of the causal chain: it should identify expected behavior, dependencies, failure modes, and the evidence needed to justify confidence.
\item Need For Cryptography and Mathematical Prelude To Cryptography address different portions of the problem. A strong comparison states their scope, assumptions, enforcement points, evidence, and how a failure in one affects the other.
\item Reviewers should identify who controls message, what inputs or dependencies it trusts, where privilege changes, what happens during partial failure, and how those assumptions are tested.
\item A weakness in mod can propagate across trust boundaries. The actual impact depends on reachable assets, granted authority, control independence, visibility, and recovery capability.
\item Useful evidence includes algorithm and mode inventories, key provenance, certificate paths, protocol traces, rotation records, and test vectors; configuration alone is insufficient when operation, monitoring, or recovery is part of the claim.
\item Prevention reduces probability, detection limits dwell time, response limits spread, and recovery restores objectives. Omitting one layer creates an unexamined failure mode.
\item Begin by establishing scope, ownership, expected behavior, and the violated assumption. Then verify the complete cryptographic construction and lifecycle rather than the algorithm name alone, preserving evidence for prioritization and follow-up.
\item Treat the exercise as stronger evidence than a static check. Prioritize according to reachable impact, likelihood of real operating conditions, missing compensating controls, and recovery difficulty.
\item Retain the principle, threat model, and control objective; separately label technologies, legal references, product examples, and threat observations whose applicability depends on time or environment.
\end{enumerate}
\section*{Key-Term Recap}
\begin{longtable}{>{\RaggedRight\arraybackslash}p{0.27\textwidth}>{\RaggedRight\arraybackslash}p{0.66\textwidth}}
\toprule\rowcolor{Navy}\color{white}\textbf{Term} & \color{white}\textbf{Meaning}\\\midrule
\endfirsthead
\toprule\rowcolor{Navy}\color{white}\textbf{Term} & \color{white}\textbf{Meaning}\\\midrule
\endhead
\textbf{Cipher} & An algorithm that transforms data using a key to provide a defined cryptographic security property. \\\midrule
\textbf{Encryption} & A keyed transformation of plaintext into ciphertext to protect confidentiality. \\\midrule
\textbf{Aes} & A standardized symmetric block cipher used to protect data when implemented with an appropriate mode and key lifecycle. \\\midrule
\textbf{Cryptography} & The use of mathematical mechanisms to protect confidentiality, integrity, authenticity, or related security properties. \\\midrule
\textbf{Authentication} & The process of establishing confidence that an asserted identity is genuine. \\\midrule
\textbf{Integrity} & The property that information and systems remain accurate, complete, and protected from unauthorized modification. \\\midrule
\textbf{Access Control} & Rules and enforcement mechanisms that restrict actions on resources to authorized subjects. \\\midrule
\textbf{Confidentiality} & The property that information is not disclosed to unauthorized parties. \\\midrule
\textbf{Digital Signature} & A public-key mechanism used to verify origin and integrity under defined key-trust assumptions. \\\midrule
\textbf{Evidence} & Observable information that supports a security conclusion and can be preserved for review. \\\midrule
\bottomrule
\end{longtable}
\begin{CornellExamBox}{One-Sentence Takeaway}
Treat data encryption as a verifiable chain from assets and assumptions through controls, evidence, response, and recovery.
\end{CornellExamBox}
\end{document}
